% GENERATED FILE. Edit canonical lessons, then run npm run generate:books.
% canonical-source-hash: fnv1a64:ed0f59fd631b2ce5
% canonical-lessons: SA-C24-water, SA-C24-rice, SA-C24-fruit, SA-S123-letter-o, SA-C24-cloth, SA-C24-lamp, SA-S122-letter-pha

\chapter{Things You Ask For}
\label{ch:sa-ask-for}
\hlchaptermodality{24}

\begin{chapteropening}
\textbf{By the end of this chapter:} \emph{I can name five everyday things I might ask someone for in Sanskrit, and say which noun family each one belongs to.}

Complete SA-C24-lamp and demonstrate the chapter promise.
\end{chapteropening}

\section[jalam]{\sk{जलम्} (jalam) — water (the stuff, not the drink)}

\label{lesson:SA-C24-water}



\begin{quote}
Before the new word: what did \emph{duhitā} mean?
\end{quote}

\subsection*{You'll want to know: \sk{जलम्}}
\textbf{\sk{जलम्}} (\emph{jalam}) — \textquotedblleft{}water (the stuff, not the drink)\textquotedblright{}.

You already have one water-word: \textbf{\sk{पानीयम्}} (\emph{pānīyam}), which names water by what you do with it — \textquotedblleft{}the thing to be drunk\textquotedblright{}. \textbf{\sk{जलम्}} names the stuff. Rain is \emph{jalam}; the river is full of \emph{jalam}; nobody offers you a cup of it in that word.

Sanskrit keeps several words here — \emph{jalam}, \emph{pānīyam}, \emph{udakam}, \emph{āpaḥ} — and a poet chooses between them the way an English poet chooses between water, drink, and the waters. The ending puts \sk{जलम्} in the neuter family you met with \sk{अन्नम्} and \sk{गृहम्}: \textbf{-\sk{अम्}}.

\begin{grammarlens}[title={What you've built}]
Two words for water now, and a reason to prefer one over the other.
\end{grammarlens}

\subsection*{Your turn}
\begin{itemize}
\raggedright
  \item \emph{Say it:} \emph{jalam}
  \item \emph{Say it:} it once more, slowly
  \item \emph{Say it:} \emph{pānīyam}, then \emph{jalam}, and say which one you would ask for at a table
  \item \emph{Recall:} say \emph{śṛṇoti}, then read \textbf{\sk{पिता}}
\end{itemize}

\begin{tcolorbox}[breakable,colback=teal!4,colframe=teal!35!black,title={Before you move on}]
What does \sk{जलम्} mean? (\textquotedblleft{}water (the stuff, not the drink)\textquotedblright{}.) Say it once more.
\end{tcolorbox}
\section[odanaḥ]{\sk{ओदनः} (odanaḥ) — cooked rice}

\label{lesson:SA-C24-rice}



\begin{quote}
Before the new word: what did \emph{jalam} name, and what did \emph{pānīyam} name?
\end{quote}

\subsection*{You'll want to know: \sk{ओदनः}}
\textbf{\sk{ओदनः}} (\emph{odanaḥ}) — \textquotedblleft{}cooked rice\textquotedblright{}.

\textbf{\sk{अन्नम्}} (\emph{annam}) was food in general — \textquotedblleft{}the eaten thing\textquotedblright{}. \textbf{\sk{ओदनः}} is narrower and warmer: the grain after it has met the water. Cooked rice, on a plate, now.

The traditional derivation takes it from a root meaning \textquotedblleft{}to wet\textquotedblright{}, which reads the word as \emph{grain that has been moistened} — a description of cooking rather than of rice. Its ending \textbf{-\sk{अः}} puts it in the masculine family beside \sk{मार्गः} and \sk{ग्रामः}, so \sk{ओदनः} and \sk{अन्नम्} sit in different families while naming the same meal.

\begin{grammarlens}[title={What you've built}]
A general word for food, and a particular thing on the plate.
\end{grammarlens}

\subsection*{Your turn}
\begin{itemize}
\raggedright
  \item \emph{Say it:} \emph{odanaḥ}
  \item \emph{Say it:} it once more, slowly
  \item \emph{Say it:} \emph{annam}, then \emph{odanaḥ}, and name the ending each one carries
  \item \emph{Recall:} read \textbf{\sk{नयति}}, then say \emph{mātā}
\end{itemize}

\begin{tcolorbox}[breakable,colback=teal!4,colframe=teal!35!black,title={Before you move on}]
What does \sk{ओदनः} mean? (\textquotedblleft{}cooked rice\textquotedblright{}.) Say it once more.
\end{tcolorbox}
\section[phalam]{\sk{फलम्} (phalam) — fruit (and the result of anything)}

\label{lesson:SA-C24-fruit}



\begin{quote}
Before the new word: what did \emph{odanaḥ} mean?
\end{quote}

\subsection*{You'll want to know: \sk{फलम्}}
\textbf{\sk{फलम्}} (\emph{phalam}) — \textquotedblleft{}fruit (and the result of anything)\textquotedblright{}.

\textbf{\sk{फलम्}} is the fruit on the tree. It is also, in the same breath and with no metaphor felt, the fruit of an action — the result, the reward, the yield. A Sanskrit sentence about the \emph{phalam} of your work is not reaching for a poetic image; it has one word where English has two.

The root behind it means \textquotedblleft{}to burst, to bear fruit\textquotedblright{}, which is the fruit split open at the moment it is ripe. Neuter \textbf{-\sk{अम्}}, like \sk{जलम्} and \sk{अन्नम्} before it.

\begin{grammarlens}[title={What you've built}]
A word that carries an orchard and a consequence at once.
\end{grammarlens}

\subsection*{Your turn}
\begin{itemize}
\raggedright
  \item \emph{Say it:} \emph{phalam}
  \item \emph{Say it:} it once more, slowly
  \item \emph{Say it:} \emph{phalam}, and then \emph{jalam} and \emph{phalam} together, naming the ending they share
  \item \emph{Recall:} say \emph{krīḍati}, then read \textbf{\sk{सूनुः}}
\end{itemize}

\begin{tcolorbox}[breakable,colback=teal!4,colframe=teal!35!black,title={Before you move on}]
What does \sk{फलम्} mean? (\textquotedblleft{}fruit (and the result of anything)\textquotedblright{}.) Say it once more.
\end{tcolorbox}
\section[o]{\sk{ओ} — one character, met inside words you already say}

\label{lesson:SA-S123-letter-o}



\begin{quote}
Before the new one, the ones you have already met: \sk{छ} · \sk{ड} · \sk{◌ू}. Say what each of them does.

One character this time. One only — and you have been saying it for pages without knowing which mark on the page it was.
\end{quote}

\subsection*{Script you'll notice: \sk{ओ}}
\textbf{\sk{ओ}} — \emph{o}.

It is an \textbf{independent vowel}: the full-size letter a vowel wears when it opens a word and has no consonant to lean on. Inside a word the same vowel is written as a small sign instead.

What it is made of:

\begin{itemize}
\raggedright
  \item the joined upper and lower bowls of \sk{अ}
  \item a separately swept middle shoulder
  \item the inner stem
  \item the trailing stem
  \item a separate upper arc
  \item the top shirorekhā
\end{itemize}

You already say these, and every one of them has \sk{ओ} somewhere inside it:

\begin{itemize}
\raggedright
  \item \textbf{\sk{ओदनः}} \emph{odanaḥ} — cooked rice
\end{itemize}

\subsection*{Writing: \sk{ओ}}
\begin{itemize}
\raggedright
  \item \textbf{1.} curve right around the upper bowl and continue down and around the lower bowl without lifting
  \item \textbf{2.} lift and sweep the middle shoulder right
  \item \textbf{3.} lift and descend the inner stem
  \item \textbf{4.} lift and descend the trailing stem
  \item \textbf{5.} lift at the trailing stem's headline junction and sweep the upper arc upward and left
  \item \textbf{6.} lift and draw the top shirorekhā left-to-right
\end{itemize}

\textbf{Pen lifts: 5.} The pen comes up 5 times and no more.

\begin{quote}
verified modern printed teaching form; traditional and everyday handwriting may divide or simplify the curves differently
\end{quote}

\begin{quote}
This is one attested teaching order and not a national standard — handwriting here is taught with school-to-school variation. Source: Saurmandal, ‘Devanagari \sk{ओ} stroke order.svg’, panels 1–6, Wikimedia Commons, 5 August 2023.
\end{quote}

\subsection*{Your turn}
\begin{itemize}
\raggedright
  \item \emph{Look:} at this, and find \sk{ओ} in it
\end{itemize}

\begin{quote}
\sk{ओदनः}
\end{quote}

\begin{itemize}
\raggedright
  \item \emph{Trace it:} \sk{ओ} three times, saying \emph{o} as you finish each one
  \item \emph{Look:} back at any page of this chapter and find \sk{ओ} once more
\end{itemize}

\begin{tcolorbox}[breakable,colback=teal!4,colframe=teal!35!black,title={Before you move on}]
Which character is this — \sk{ओ}? What sound does it carry? (\emph{\textbf{o}}.) Name one word you already say that contains it.
\end{tcolorbox}
\section[vastram]{\sk{वस्त्रम्} (vastram) — cloth, a garment}

\label{lesson:SA-C24-cloth}



\begin{quote}
Before the new word: what did \emph{phalam} mean — both of its meanings?
\end{quote}

\subsection*{You'll want to know: \sk{वस्त्रम्}}
\textbf{\sk{वस्त्रम्}} (\emph{vastram}) — \textquotedblleft{}cloth, a garment\textquotedblright{}.

You know \textbf{\sk{वसामि}} (\emph{vasāmi}), \textquotedblleft{}I dwell\textquotedblright{}, from the root \emph{vas}. \textbf{\sk{वस्त्रम्}} looks like it belongs to that root and does not. Sanskrit has two roots spelled \emph{vas}: one for dwelling, one for wearing. \sk{वसामि} belongs to the first, \sk{वस्त्रम्} to the second, and no amount of staring at the letters will separate them — only the meanings do.

The wearing one travelled. Latin \emph{vestis}, and through it English \emph{vest}, \emph{vestment}, \emph{invest} — which originally meant to clothe someone in the robes of an office. Neuter \textbf{-\sk{अम्}}.

\begin{grammarlens}[title={What you've built}]
A word that warns you: same letters, different root.
\end{grammarlens}

\subsection*{Your turn}
\begin{itemize}
\raggedright
  \item \emph{Say it:} \emph{vastram}
  \item \emph{Say it:} it once more, slowly
  \item \emph{Say it:} \emph{vasāmi}, then \emph{vastram}, and say which root each one belongs to
  \item \emph{Recall:} read \textbf{\sk{दहति}}, then say \emph{duhitā}
\end{itemize}

\begin{tcolorbox}[breakable,colback=teal!4,colframe=teal!35!black,title={Before you move on}]
What does \sk{वस्त्रम्} mean? (\textquotedblleft{}cloth, a garment\textquotedblright{}.) Say it once more.
\end{tcolorbox}
\section[dīpaḥ]{\sk{दीपः} (dīpaḥ) — a lamp}

\label{lesson:SA-C24-lamp}



\begin{quote}
Before the last word of this run: what did \emph{vastram} mean, and what did \emph{odanaḥ} mean?
\end{quote}

\subsection*{You'll want to know: \sk{दीपः}}
\textbf{\sk{दीपः}} (\emph{dīpaḥ}) — \textquotedblleft{}a lamp\textquotedblright{}.

\textbf{\sk{दीपः}} is a lamp — the small clay one with oil and a wick. Its root means \textquotedblleft{}to blaze\textquotedblright{}, so the word names the burning rather than the vessel.

Put \sk{दीप} together with a word meaning \textquotedblleft{}a row\textquotedblright{} and you get \emph{dīpāvalī}, \textquotedblleft{}a row of lamps\textquotedblright{} — the festival English writes as Diwali. Masculine \textbf{-\sk{अः}}, like \sk{ओदनः}.

\begin{grammarlens}[title={What you've built}]
Five things you might ask for: water, rice, fruit, cloth, a lamp.
\end{grammarlens}

\subsection*{Your turn}
Say these aloud:

\begin{itemize}
\raggedright
  \item \emph{dīpaḥ}
  \item it once more, slowly
  \item all five in order, and sort them into the two noun families
\end{itemize}

\begin{tcolorbox}[breakable,colback=teal!4,colframe=teal!35!black,title={Before you move on}]
What does \sk{दीपः} mean? (\textquotedblleft{}a lamp\textquotedblright{}.) Say it once more.
\end{tcolorbox}
\section[pha]{\sk{फ} — one character, met inside words you already say}

\label{lesson:SA-S122-letter-pha}



\begin{quote}
Before the new one, the ones you have already met: \sk{ड} · \sk{◌ू} · \sk{ओ}. Say what each of them does.

One character this time. One only — and you have been saying it for pages without knowing which mark on the page it was.
\end{quote}

\subsection*{Script you'll notice: \sk{फ}}
\textbf{\sk{फ}} — \emph{pha}.

It is a \textbf{consonant}, and in this script a consonant is never bare: it comes with an \emph{a} already in it. So it is not \emph{ph}, it is \textbf{pha}.

What it is made of:

\begin{itemize}
\raggedright
  \item a left stem joined through the lower bowl to a rising and retraced central stem
  \item a separate clockwise right arch
  \item the top shirorekhā
\end{itemize}

You already say these, and every one of them has \sk{फ} somewhere inside it:

\begin{itemize}
\raggedright
  \item \textbf{\sk{फलम्}} \emph{phalam} — fruit
\end{itemize}

\subsection*{Writing: \sk{फ}}
\begin{itemize}
\raggedright
  \item \textbf{1.} descend the left stem, curve right around the lower bowl, rise along the central side to the headline, then retrace down the central stem without lifting
  \item \textbf{2.} lift at the central junction and sweep clockwise over and down through the open right arch
  \item \textbf{3.} lift and draw the top shirorekhā left-to-right
\end{itemize}

\textbf{Pen lifts: 2.} The pen comes up 2 times and no more.

\begin{quote}
verified three-stroke teaching form; another learner source stages the components differently
\end{quote}

\begin{quote}
This is one attested teaching order and not a national standard — handwriting here is taught with school-to-school variation. Source: JackPotte, ‘Devanagari p\textsuperscript{h} \sk{फ}.gif’, strokes 1–3, Wikimedia Commons, 29 March 2009.
\end{quote}

\subsection*{Your turn}
\begin{itemize}
\raggedright
  \item \emph{Look:} at this, and find \sk{फ} in it
\end{itemize}

\begin{quote}
\sk{फलम्}
\end{quote}

\begin{itemize}
\raggedright
  \item \emph{Trace it:} \sk{फ} three times, saying \emph{pha} as you finish each one
  \item \emph{Look:} back at any page of this chapter and find \sk{फ} once more
\end{itemize}

\begin{tcolorbox}[breakable,colback=teal!4,colframe=teal!35!black,title={Before you move on}]
Which character is this — \sk{फ}? What sound does it carry? (\emph{\textbf{pha}}.) Name one word you already say that contains it.
\end{tcolorbox}
